% ===============================================================
%  PART VII -- Limitations and Conclusion (Chapters 21--22)
% ===============================================================

\part{Limitations and Conclusion}
\label{part:conclusion}

\begin{center}
\adjustbox{max width=\linewidth}{%
\begin{forest}
  part map,
  [Part VII
    [Ch.~21: Limitations
      [Modelled vs left out] [Threats to validity] [Claims not made]]
    [Ch.~22: Conclusion
      [Answers to RQ1--RQ4] [Contributions] [Future work]]
  ]
\end{forest}}
\end{center}

\begin{keybox}[What Part VII does]
Chapter~21 sets out what the study cannot show: what the simulator leaves out,
which threats could make a result wrong and what we did about each, and which
claims we do not make. Chapter~22 answers the four research questions of
Chapter~10 and lists the work that would test the conclusions further.
\end{keybox}

% ===============================================================
%  Chapter 21 -- Limitations and Threats to Validity
% ===============================================================
\chapter{Limitations and Threats to Validity}
\label{ch:limits}

\begin{figure}[htbp]
\centering
\adjustbox{max width=\linewidth}{%
\begin{forest}
  chapter map,
  [Chapter 21
    [Modelled vs left out [What the simulator has] [What it lacks]]
    [Threats to validity [Internal] [Construct] [External]]
    [Claims not made [Absolute times] [Named programs] [Universal DSR line]]
  ]
\end{forest}}
\caption[]{Structure of Chapter~21: what the model contains,
what could make its results wrong, and the limits on what can be concluded.}
\label{fig:ch21-map}
\end{figure}

% ---------------------------------------------------------------
\section{What the Study Models and What It Leaves Out}
\label{sec:lim-model}
\sectionmap{[Model [{Private caches, one bus, one memory port}]
  [{No pipeline, no DRAM, no directory}] [{Synthetic workloads}]
  [Effect on the results]]}
\index{MOESI}\index{directory protocol}%

\Cref{fig:lim-model} sets what the two layers of Chapter~17 contain beside what a
real multicore processor also has. Some omissions would change the size of the
results in a predictable direction. A core that keeps working past a miss, through
out-of-order execution or prefetching, hides part of every miss latency, so the
small time gains measured at low load would probably shrink further; at saturation,
where time is set by bus occupancy, the limit of \cref{eq:sat-limit} would still
apply. A shared LLC would make memory-side reads faster and less congested, which
by experiments E4 and E5 would reduce MOESI's time advantage, the trend Fagundes
and Souza report~\cite{fagundes2024evaluation}. Other omissions, such as real
programs and directories, change what is measured rather than its size.

\begin{figure}[htbp]
\centering
\begin{bigfig}[0.95]
\begin{tikzpicture}[font=\sffamily\scriptsize,
    item/.style={draw, thick, rounded corners=2pt, minimum width=5.4cm, minimum height=0.55cm, align=left, anchor=west, inner sep=4pt}]
  \node[font=\sffamily\bfseries\small, text=memGreen!80!black] at (2.7,5.0) {Modelled};
  \node[font=\sffamily\bfseries\small, text=alertRed] at (9.6,5.0) {Left out};
  \foreach \y/\t in {4.3/{private write-back caches, LRU, 64-B blocks},
                     3.6/{MSI, MESI, MOESI differing in three rules},
                     2.9/{one snooping bus; messages priced in bytes},
                     2.2/{one memory port; owner cache ports},
                     1.5/{per-core clocks; queueing for shared resources},
                     0.8/{report latencies $T_{L1}$, $T_{\mathrm{mem}}$, $T_{\mathrm{coh}}$},
                     0.1/{synthetic mixtures of four sharing patterns}} {
    \node[item, fill=memGreen!10, draw=memGreen] at (0,\y) {\t};
  }
  \foreach \y/\t in {4.3/{shared last-level cache; inclusion},
                     3.6/{processor pipeline, out-of-order execution},
                     2.9/{prefetching; DRAM banks and row buffers},
                     2.2/{directories and point-to-point networks},
                     1.5/{other MOESI members (forwarding, MOESI-prime)},
                     0.8/{false sharing; races between requests},
                     0.1/{real programs, operating system, energy}} {
    \node[item, fill=alertRed!8, draw=alertRed] at (6.9,\y) {\t};
  }
  \draw[thick, coreGrey, dashed] (6.45,-0.35) -- (6.45,4.7);
\end{tikzpicture}
\end{bigfig}
\caption[What the study models]{What the simulator of Part~V contains (left) and
what a real multicore processor also has (right). Each item on the right is a
reason why Part~VI describes trends, not the behaviour of a particular machine.}
\label{fig:lim-model}
\end{figure}

% ---------------------------------------------------------------
\section{Threats to Validity}
\label{sec:lim-threats}
\sectionmap{[Threats [{Internal: is the result right for the model?}]
  [{Construct: is the right thing measured?}] [{External: does it carry over?}]
  [Mitigation for each]]}
\index{forced writeback}\index{Dirty Sharing Ratio}%

\Cref{tab:threats} lists the threats we identified, in the usual three groups:
internal (could the result be wrong even for the model?), construct (does the
metric measure what we claim?) and external (would the result hold outside the
model?), with what we did about each. The external threat with the largest
consequence is the first one. Every number in Part~VI comes from synthetic traces
whose sharing we chose. They contain the patterns that, by Chapter~9, decide the
cost of the forced writeback, and DSR was tested over a wide range of mixes, core
counts and cache sizes, but a real program has working sets, phases and
synchronisation that no mix of four patterns reproduces. The DSR line was fitted and
tested on data from one simulator, and its slope and accuracy on real programs are
unknown.

\begin{table}[htbp]
\centering
\caption{Threats to the validity of the results, and what was done about each.}
\label{tab:threats}
\footnotesize
\begin{tabular}{@{}L{1.5cm} L{4.6cm} L{6.6cm}@{}}
\toprule
Kind & Threat & Mitigation \\
\midrule
Internal & a bug in the simulator could look like a protocol effect &
  22 analytical checks, coherence checked by version numbers on every access,
  the Part~IV identities on 468 pairs, 11 timing checks (\cref{sec:sim-verify}) \\
 & in the timing study the two protocols' interleavings can differ (read
  misses differ by up to 1.8\,\%) & exact claims made for the traffic study only;
  timing reported as paired ratios over seeds (\cref{sec:method-timing}) \\
 & only three seeds per configuration & errors on held-out seeds match in-sample
  errors (\cref{sec:res-h3}); spreads reported, not only means \\
 & the write-miss transfer rule was chosen by us & identical in all protocols;
  not exercised by the validated patterns (\cref{sec:sim-traffic}) \\
\midrule
Construct & $F$ and the saving agree almost by construction &
  stated explicitly (\cref{sec:res-h3}); the DSR line, which uses only
  transaction counts, is the result tested \\
 & execution time is the finish time of the slowest core & comparisons are
  paired on identical streams, so imbalance affects both protocols \\
 & bytes are not energy, and AMAT is not time & time measured separately in the
  timing layer; no energy claims; AMAT and time compared in \cref{sec:res-amat} \\
\midrule
External & synthetic workloads, not real programs & built from the four
  patterns that decide the cost (Chapter~9); results stated per pattern mix \\
 & no pipeline, prefetching, DRAM or shared last-level cache & results stated
  as trends and relative differences; expected direction of each omission
  discussed in \cref{sec:lim-model} \\
 & assumed bus, upgrade and port latencies & each swept: E1 (load), E4
  (cache-to-cache, owner port), E5 (memory port) \\
 & one snooping bus; no directory & directory behaviour left to the analysis
  of Chapter~13 and \cref{tab:p2p} \\
 & one member of the MOESI class & chosen deliberately (\cref{sec:method-delta});
  others can be studied by the same method \\
 & 2 to 16 cores only & the range set in Chapter~10; trends beyond 16 not
  claimed \\
\bottomrule
\end{tabular}
\end{table}

% ---------------------------------------------------------------
\section{Claims This Report Does Not Make}
\label{sec:lim-claims}
\sectionmap{[Not claimed [Absolute execution times] [Results for named programs]
  [{MOESI always faster}] [{One universal DSR line}]]}
\index{MESI}\index{MOESI}\index{Dirty Sharing Ratio}\index{AMAT}\index{directory protocol}\index{false sharing}\index{memory consistency}%

We do not claim that any real processor would show the execution times, speed-ups
or AMAT values of Part~VI; only their trends and relative sizes are claimed. We do
not claim that any named application or benchmark has a particular DSR or would
gain a particular amount from MOESI. MOESI is not always faster than MESI, as E4
shows (\cref{sec:res-h4}). The fitted line $S_B \approx 92.6\,\dsr - 3.6$ is not
universal, since its slope depends on $\bar b$ and therefore on the system
(\cref{sec:res-h3}). And we make no claims about energy, directories, consistency
models or false sharing, which lie outside the study (\cref{tab:scope-change}).

% ===============================================================
%  Chapter 22 -- Conclusion and Future Work
% ===============================================================
\chapter{Conclusion and Future Work}
\label{ch:conclusion}

\begin{figure}[htbp]
\centering
\adjustbox{max width=\linewidth}{%
\begin{forest}
  chapter map,
  [Chapter 22
    [Answers [RQ1: how large] [RQ2: does it become time] [RQ3: can it be predicted] [RQ4: when does it hurt]]
    [Contributions [Analytical] [Method {\&} metric] [Simulator]]
    [Future work [Real traces] [gem5 cross-check] [Directories {\&} variants]]
  ]
\end{forest}}
\caption[]{Structure of Chapter~22: answers to the four
research questions, the contributions, and the next steps.}
\label{fig:ch22-map}
\end{figure}

% ---------------------------------------------------------------
\section{Answers to the Research Questions}
\label{sec:conc-problem}
\label{sec:conc-answers}
\sectionmap{[Answers [{RQ1: a third to a half, falling with cores}]
  [RQ2: only under load] [{RQ3: yes, through DSR}] [RQ4: only if c2c is slow]]}
\index{MESI}\index{MOESI}\index{forced writeback}\index{dirty sharing}\index{Dirty Sharing Ratio}\index{survival fraction}\index{dirty epoch}\index{coherence axioms}\index{cache-to-cache transfer}\index{owner}%

The report started from one transaction. When a core reads a block that another
core holds in \stM, MESI supplies the block and also writes it back to memory,
because its shared state must be clean (\cref{fig:msi}). The coherence axioms do not
require this writeback (Chapter~10), and the owner-based protocols of Part~III
remove it by keeping one cache responsible for the dirty data.
\Cref{tab:conc-answers} gives the four answers.

\begin{table}[htbp]
\centering
\caption{The research questions of \cref{tab:rq} and their answers.}
\label{tab:conc-answers}
\small
\begin{tabular}{@{}L{1.0cm} L{4.0cm} L{6.0cm} L{1.9cm}@{}}
\toprule
 & Question & Answer & Evidence \\
\midrule
RQ1 & How large is $F(W,N)$, and how does it change from 2 to 16 cores? &
  33--44\,\% of MESI's coherence bytes when dirty sharing is substantial;
  falls as cores are added & \cref{sec:res-h1} \\
RQ2 & Does saved traffic become saved time? & only when the bus or memory is
  busy: about 1\,\% of the saving at low utilisation, 95\,\% at saturation &
  \cref{sec:res-h2,sec:res-memory} \\
RQ3 & Can the benefit be predicted from MESI behaviour? & yes for traffic:
  exactly by $\sigma\cdot72\,\dsr/\bar b$, and to 1.8 points by a line in DSR
  alone & \cref{sec:res-h3} \\
RQ4 & Can deferring the writeback cost more than it saves? & never in bytes;
  in time only if a cache-to-cache read is much slower than a memory read &
  \cref{sec:res-sigma,sec:res-h4} \\
\bottomrule
\end{tabular}
\end{table}

When dirty sharing is a large part of a workload, the forced writeback is a third
to nearly half of the bytes MESI's protocol moves, 43.6\,\% on average at 2 cores
and 33.4\,\% at 16. Its share falls with more cores because the first reader in each
epoch pays for the writeback and later readers find a clean block. It is also nearly
all of MESI's memory writes in these workloads, which is why MOESI removed 83 to
99.8\,\% of them.

Removing it saves time only when a shared resource is busy. MOESI removed about
30\,\% of the bytes at every load but cut execution time by 0.3\,\% on an idle bus
and by 28\,\% on a saturated one, where the speed-up reached the bus-occupancy limit
derived in Part~IV. On a system limited by memory bandwidth the gain was 36\,\%. A
comparison on a lightly loaded system will see almost no difference in time, as the
MOESI-prime evaluation did~\cite{loughlin2022moesiprime}.

The traffic benefit can be predicted. Under the controlled method the two protocols
hold the same blocks at every step and the saving equals
$\sigma\cdot72\,\dsr/\bar b$ exactly. The Dirty Sharing Ratio alone, counted with
two counters on a MESI system, predicted the saving through one straight line with
a mean error of 1.84 points and kept that accuracy on core counts, traces and cache
sizes it was not fitted to. The share of operations touching shared data does not
work as a predictor. Time cannot be predicted from DSR alone, since it also depends
on how busy the bus or memory is (\cref{fig:res-rule}).

MOESI cannot send more bytes than MESI under the method of Chapter~15, and it did
not in any of 756 paired runs. It was slower only when a cache-to-cache read was
made much slower than a memory read, above about 140 cycles against 100, because
the owner then slowly answers reads that MESI would have served from memory. A
measurement of the dirty reads alone would not show this loss. Small caches reduce
MOESI's gain without reversing it, mainly by lowering the number of dirty reads.

% ---------------------------------------------------------------
\section{What This Report Contributes}
\label{sec:conc-contrib}
\sectionmap{[Contributions [Analytical results in Parts I--III]
  [{Method, proofs and metric in Part IV}] [Verified simulator in Part V]
  [Tested hypotheses in Part VI]]}
\index{MESI}\index{MOESI}\index{Owned state}\index{verification}\index{proof}%

\Cref{tab:contrib-list} lists the contributions. The report completes the MESI
protocol that the textbook leaves as an exercise and prices every protocol reviewed
on one set of counting rules (C1--C10). It proposes a controlled method in which
MOESI differs from MESI in three rules, and proves that the two protocols then hold
the same blocks and differ only in writebacks (C11, C12). It introduces the Dirty
Sharing Ratio, proves that MOESI never sends more bytes than MESI and derives the
exact saving identity (C13, C14). It separates the Owned state into two effects with
different signs, which explains when MOESI can be slower (C15). And it builds and
verifies a two-layer simulator and uses it to test four hypotheses stated before
any measurement (C16). The proofs are counting arguments and the simulator is a few
hundred lines whose rules are the tables of this report, so all of it is within the
reach of an undergraduate computer architecture course.

% ---------------------------------------------------------------
\section{Future Work}
\label{sec:conc-future}
\label{sec:conc-close}
\sectionmap{[Future work [Real program traces] [Cycle-accurate cross-check]
  [Directories and MOESI variants] [Energy and false sharing]]}
\index{MESI}\index{MOESI}\index{MESIF}\index{Owned state}\index{dirty sharing}\index{Dirty Sharing Ratio}\index{directory protocol}\index{false sharing}\index{OpenMP}\index{gem5}\index{owner}%

The limitations of Chapter~21 set the next steps (\cref{fig:conc-future}). The most
important is to record memory traces of the OpenMP programs planned in Chapter~10
with a binary-instrumentation tool, replay them through both layers, and test
whether the DSR line keeps its accuracy and slope. Next, the three changed rules of
\cref{tab:transition-diff} could be built into one protocol of a detailed simulator
such as \texttt{gem5}~\cite{binkert2011gem5}, so the controlled comparison can be
repeated with pipelines, a shared LLC and DRAM. The method could then be applied to
a directory protocol, where the dirty read takes three hops
(\cref{ch:scaling-approaches}), and to other members of the owner family, adding one
feature at a time, such as the clean forwarding of MESIF~\cite{goodman2009mesif} or
the extra states of MOESI-prime~\cite{loughlin2022moesiprime}. Finally, each protocol
could be priced in energy as well as bytes, and false sharing, left out here, could
be examined for its effect on DSR.

\begin{figure}[htbp]
\centering
\adjustbox{max width=0.9\linewidth}{%
\begin{forest}
  chapter map,
  [Future work
    [Real workloads
      [{Record traces of OpenMP programs}]
      [{Test the DSR line on them}]]
    [Detailed timing
      [{Cross-check trends in gem5}]
      [{Add pipeline, LLC, DRAM}]]
    [Wider design space
      [{Directory and token versions}]
      [{MESIF, MOESI-prime, forwarding}]]
    [New questions
      [{Energy per protocol}]
      [{False sharing and DSR}]]
  ]
\end{forest}}
\caption[Future work]{Directions for future work, grouped by the limitation of
Chapter~21 each one addresses.}
\label{fig:conc-future}
\end{figure}

In short, the Owned state adds one state and changes three rules, and it does two
separate things. It removes a writeback that coherence never required, which always
reduces traffic and memory writes and saves time when bandwidth is scarce. It also
makes the owner, instead of memory, answer later readers, which helps or hurts
depending on which of the two is faster. Which protocol is better therefore depends
on the workload's dirty sharing, the system's load and its latencies, and the first
of these can be measured on a MESI system, before MOESI is built, by counting two
events.
